\documentclass[a4paper,12pt,twoside]{memoir}

% Castellano
\usepackage[spanish,es-tabla]{babel}
\selectlanguage{spanish}
\usepackage[utf8]{inputenc}
\usepackage[T1]{fontenc}
\usepackage{lmodern} % Scalable font
\usepackage{microtype}
\usepackage{placeins}

\RequirePackage{booktabs}
\RequirePackage[table]{xcolor}
\RequirePackage{xtab}
\RequirePackage{multirow}

% Links
\usepackage[colorlinks]{hyperref}
\hypersetup{
	allcolors = {red}
}

% Ecuaciones
\usepackage{amsmath}

% Rutas de fichero / paquete
\newcommand{\ruta}[1]{{\sffamily #1}}

% Párrafos
\nonzeroparskip


% Imagenes
\usepackage{graphicx}
\newcommand{\imagen}[2]{
	\begin{figure}[!h]
		\centering
		\includegraphics[width=0.9\textwidth]{#1}
		\caption{#2}\label{fig:#1}
	\end{figure}
	\FloatBarrier
}

\newcommand{\imagenflotante}[2]{
	\begin{figure}%[!h]
		\centering
		\includegraphics[width=0.9\textwidth]{#1}
		\caption{#2}\label{fig:#1}
	\end{figure}
}



% El comando \figura nos permite insertar figuras comodamente, y utilizando
% siempre el mismo formato. Los parametros son:
% 1 -> Porcentaje del ancho de página que ocupará la figura (de 0 a 1)
% 2 --> Fichero de la imagen
% 3 --> Texto a pie de imagen
% 4 --> Etiqueta (label) para referencias
% 5 --> Opciones que queramos pasarle al \includegraphics
% 6 --> Opciones de posicionamiento a pasarle a \begin{figure}
\newcommand{\figuraConPosicion}[6]{%
  \setlength{\anchoFloat}{#1\textwidth}%
  \addtolength{\anchoFloat}{-4\fboxsep}%
  \setlength{\anchoFigura}{\anchoFloat}%
  \begin{figure}[#6]
    \begin{center}%
      \Ovalbox{%
        \begin{minipage}{\anchoFloat}%
          \begin{center}%
            \includegraphics[width=\anchoFigura,#5]{#2}%
            \caption{#3}%
            \label{#4}%
          \end{center}%
        \end{minipage}
      }%
    \end{center}%
  \end{figure}%
}

%
% Comando para incluir imágenes en formato apaisado (sin marco).
\newcommand{\figuraApaisadaSinMarco}[5]{%
  \begin{figure}%
    \begin{center}%
    \includegraphics[angle=90,height=#1\textheight,#5]{#2}%
    \caption{#3}%
    \label{#4}%
    \end{center}%
  \end{figure}%
}
% Para las tablas
\newcommand{\otoprule}{\midrule [\heavyrulewidth]}
%
% Nuevo comando para tablas pequeñas (menos de una página).
\newcommand{\tablaSmall}[5]{%
 \begin{table}
  \begin{center}
   \rowcolors {2}{gray!35}{}
   \begin{tabular}{#2}
    \toprule
    #4
    \otoprule
    #5
    \bottomrule
   \end{tabular}
   \caption{#1}
   \label{tabla:#3}
  \end{center}
 \end{table}
}

%
% Nuevo comando para tablas pequeñas (menos de una página).
\newcommand{\tablaSmallSinColores}[5]{%
 \begin{table}[H]
  \begin{center}
   \begin{tabular}{#2}
    \toprule
    #4
    \otoprule
    #5
    \bottomrule
   \end{tabular}
   \caption{#1}
   \label{tabla:#3}
  \end{center}
 \end{table}
}

\newcommand{\tablaApaisadaSmall}[5]{%
\begin{landscape}
  \begin{table}
   \begin{center}
    \rowcolors {2}{gray!35}{}
    \begin{tabular}{#2}
     \toprule
     #4
     \otoprule
     #5
     \bottomrule
    \end{tabular}
    \caption{#1}
    \label{tabla:#3}
   \end{center}
  \end{table}
\end{landscape}
}

%
% Nuevo comando para tablas grandes con cabecera y filas alternas coloreadas en gris.
\newcommand{\tabla}[6]{%
  \begin{center}
    \tablefirsthead{
      \toprule
      #5
      \otoprule
    }
    \tablehead{
      \multicolumn{#3}{l}{\small\sl continúa desde la página anterior}\\
      \toprule
      #5
      \otoprule
    }
    \tabletail{
      \hline
      \multicolumn{#3}{r}{\small\sl continúa en la página siguiente}\\
    }
    \tablelasttail{
      \hline
    }
    \bottomcaption{#1}
    \rowcolors {2}{gray!35}{}
    \begin{xtabular}{#2}
      #6
      \bottomrule
    \end{xtabular}
    \label{tabla:#4}
  \end{center}
}

%
% Nuevo comando para tablas grandes con cabecera.
\newcommand{\tablaSinColores}[6]{%
  \begin{center}
    \tablefirsthead{
      \toprule
      #5
      \otoprule
    }
    \tablehead{
      \multicolumn{#3}{l}{\small\sl continúa desde la página anterior}\\
      \toprule
      #5
      \otoprule
    }
    \tabletail{
      \hline
      \multicolumn{#3}{r}{\small\sl continúa en la página siguiente}\\
    }
    \tablelasttail{
      \hline
    }
    \bottomcaption{#1}
    \begin{xtabular}{#2}
      #6
      \bottomrule
    \end{xtabular}
    \label{tabla:#4}
  \end{center}
}

%
% Nuevo comando para tablas grandes sin cabecera.
\newcommand{\tablaSinCabecera}[5]{%
  \begin{center}
    \tablefirsthead{
      \toprule
    }
    \tablehead{
      \multicolumn{#3}{l}{\small\sl continúa desde la página anterior}\\
      \hline
    }
    \tabletail{
      \hline
      \multicolumn{#3}{r}{\small\sl continúa en la página siguiente}\\
    }
    \tablelasttail{
      \hline
    }
    \bottomcaption{#1}
  \begin{xtabular}{#2}
    #5
   \bottomrule
  \end{xtabular}
  \label{tabla:#4}
  \end{center}
}



\definecolor{cgoLight}{HTML}{EEEEEE}
\definecolor{cgoExtralight}{HTML}{FFFFFF}

%
% Nuevo comando para tablas grandes sin cabecera.
\newcommand{\tablaSinCabeceraConBandas}[5]{%
  \begin{center}
    \tablefirsthead{
      \toprule
    }
    \tablehead{
      \multicolumn{#3}{l}{\small\sl continúa desde la página anterior}\\
      \hline
    }
    \tabletail{
      \hline
      \multicolumn{#3}{r}{\small\sl continúa en la página siguiente}\\
    }
    \tablelasttail{
      \hline
    }
    \bottomcaption{#1}
    \rowcolors[]{1}{cgoExtralight}{cgoLight}

  \begin{xtabular}{#2}
    #5
   \bottomrule
  \end{xtabular}
  \label{tabla:#4}
  \end{center}
}


















\graphicspath{ {./img/} }

% Capítulos
\chapterstyle{bianchi}
\newcommand{\capitulo}[2]{
	\setcounter{chapter}{#1}
	\setcounter{section}{0}
	\chapter*{#2}
	\addcontentsline{toc}{chapter}{#2}
	\markboth{#2}{#2}
}

% Apéndices
\renewcommand{\appendixname}{Apéndice}
\renewcommand*\cftappendixname{\appendixname}

\newcommand{\apendice}[1]{
	%\renewcommand{\thechapter}{A}
	\chapter{#1}
}

\renewcommand*\cftappendixname{\appendixname\ }

% Formato de portada
\makeatletter
\usepackage{xcolor}
\newcommand{\tutor}[1]{\def\@tutor{#1}}
\newcommand{\course}[1]{\def\@course{#1}}
\definecolor{cpardoBox}{HTML}{E6E6FF}
\def\maketitle{
  \null
  \thispagestyle{empty}
  % Cabecera ----------------
\noindent\includegraphics[width=\textwidth]{cabecera}\vspace{1cm}%
  \vfill
  % Título proyecto y escudo informática ----------------
  \colorbox{cpardoBox}{%
    \begin{minipage}{.8\textwidth}
      \vspace{.5cm}\Large
      \begin{center}
      \textbf{TFG del Grado en Ingeniería Informática}\vspace{.6cm}\\
      \textbf{\LARGE\@title{}}
      \end{center}
      \vspace{.2cm}
    \end{minipage}

  }%
  \hfill\begin{minipage}{.20\textwidth}
    \includegraphics[width=\textwidth]{escudoInfor}
  \end{minipage}
  \vfill
  % Datos de alumno, curso y tutores ------------------
  \begin{center}%
  {%
    \noindent\LARGE
    Presentado por \@author{}\\ 
    en Universidad de Burgos --- \@date{}\\
    Tutor: \@tutor{}\\
  }%
  \end{center}%
  \null
  \cleardoublepage
  }
\makeatother

\newcommand{\nombre}{Jesús Manuel Calvo Ruiz de Temiño} %%% cambio de comando
\newcommand{\titulo}{Ampliación, actualización y mantenimiento de la aplicación \textsl{TourPlanner}. Parte FrontEnd}
\newcommand{\mtutor}{Bruno Baruque Zanon}

% Datos de portada
\title{\titulo}
\author{\nombre}
\tutor{\mtutor}
\date{\today}

\begin{document}

\maketitle


\newpage\null\thispagestyle{empty}\newpage


%%%%%%%%%%%%%%%%%%%%%%%%%%%%%%%%%%%%%%%%%%%%%%%%%%%%%%%%%%%%%%%%%%%%%%%%%%%%%%%%%%%%%%%%
\thispagestyle{empty}


\noindent\includegraphics[width=\textwidth]{cabecera}\vspace{1cm}

\noindent D. \mtutor, profesor del departamento de nombre departamento, área de nombre área.

\noindent Expone:

\noindent Que el alumno D. \nombre, con DNI 13173674X, ha realizado el Trabajo final de Grado en Ingeniería Informática titulado \titulo. 

\noindent Y que dicho trabajo ha sido realizado por el alumno bajo la dirección del que suscribe, en virtud de lo cual se autoriza su presentación y defensa.

\begin{center} %\large
En Burgos, {\large \today}
\end{center}

\vfill\vfill\vfill

% Author and supervisor
%\begin{minipage}{0.45\textwidth}
%\begin{flushleft} %\large
%Vº. Bº. del Tutor:\\[2cm]
%D. nombre tutor
%\end{flushleft}
%\end{minipage}
%\hfill
%\begin{minipage}{0.45\textwidth}
%\begin{flushleft} %\large
%Vº. Bº. del co-tutor:\\[2cm]
%. nombre co-tutor
%\end{flushleft}
%\end{minipage}
%\hfill

\vfill

% para casos con solo un tutor comentar lo anterior
% y descomentar lo siguiente
Vº. Bº. del Tutor:\\[2cm]
D. nombre tutor


\newpage\null\thispagestyle{empty}\newpage




\frontmatter

% Abstract en castellano
\renewcommand*\abstractname{Resumen}
\begin{abstract}
Este proyecto se basa en realizar la ampliación, actualización y mantenimiento de la aplicación TourPlanner, la cual se encarga de la generación de rutas turísticas personalizadas a través de una serie de algoritmos de búsqueda de caminos óptimos y que esta diseñada para ser utilizada en sistemas operativos Android.

Para desarrollar este trabajo hemos cooperado mi compañero, Ignacio Aparicio Blanco y yo, pero cada uno realizando dos partes separadas que finalmente cumplen una función juntas, ya que Ignacio ha realizado la parte del backEnd (algoritmos para la generación de las rutas y manejo de la base de datos) y yo me he encargado de la parte frontEnd (interfaz y funcionamiento interno de la aplicación Android).

A continuación explicaré brevemente las mejoras que he realizado en la parte FrontEnd de la aplicación. Cabe destacar que este trabajo esta centrado en la aplicación de mejoras estructurales y en la actualización y estandarización de las librerías y métodos utilizados, para así conseguir un mantenimiento futuro mucho mas sencillo.

Se ha conseguido actualizar la aplicación para que pueda funcionar en los dispositivos mas modernos, ya que esta desarrollada para la versión de Android 6.0, pero también es compatible hasta la versión Android 9.

También se ha hecho uso de una herramienta estándar para la creación del menú de la aplicación, que se encuentra dentro del paquete de Android. Con esto se consigue que la aplicación tenga una apariencia mucho más parecida a otras aplicaciones muy utilizadas en sistemas Android, como puede ser el propio Google Maps o la aplicación de Gmail, lo cual hace que resulte mas fácil de entender su funcionamiento para el usuario medio.

Otra parte de este proyecto es bastante centrada en la optimización con respecto a los recursos y al uso de ciertos métodos que contiene Android. Se ha cambiado la estructura de "activities" dentro de la aplicación, para pasar a usar una estructura mas óptima.Más adelante se explicará en que consiste.

Por otro lado, como ya he dicho, este trabajo se centra en la estandarización, y por ello se han dejado de utilizar una serie de librerías que, ya sea por no tener continuidad o por no tener compatibilidad con las nuevas versiones de Android, no eran capaces de funcionar en este sistema y por tanto, se ha decidido buscar la opción mas estándar o sostenible para sustituirlas.


\end{abstract}

\renewcommand*\abstractname{Descriptores}
\begin{abstract}
Android, actualización, mantenimiento, optimización, menú, estandarización, frontEnd

\ldots
\end{abstract}

\clearpage

% Abstract en inglés
\renewcommand*\abstractname{Abstract}
\begin{abstract}
A \textbf{brief} presentation of the topic addressed in the project.
\end{abstract}

\renewcommand*\abstractname{Keywords}
\begin{abstract}
keywords separated by commas.
\end{abstract}

\clearpage

% Indices
\tableofcontents

\clearpage

\listoffigures

\clearpage

\listoftables
\clearpage

\mainmatter
\capitulo{1}{Introducción}

La aplicación sobre la que se desarrollará este proyecto de fin de grado comenzó con el proyecto \textit{Generación de rutas Turísticas personalizadas} \cite{tfg1} y un año después fue ampliado en el trabajo \textit{Ampliación sobre la aplicación para la generación de rutas turísticas personalizadas} \cite{tfm1}. En este caso se trabajará sobre la última versión de la que se dispone, donde se realizarán todas las mejoras y actualizaciones necesarias, ya que se trata de un proyecto realizado sobre tecnologías con varios años de antigüedad.

Una parte del proyecto se basa en la generación de rutas turísticas óptimas respecto a los gustos del usuario, añadiendo el factor de horarios de apertura y cierre de los puntos de interés. Esto se consigue con el desarrollo de una serie de algoritmos que se han desarrollado en el trabajo de BackEnd \cite{tfg2}.

Por el otro lado, la interfaz completa de la aplicación, donde el usuario podrá realizar toda la interacción con el sistema sería la correspondiente a la aplicación de Android. Se ha desarrollado de forma que cualquier usuario que tenga costumbre en el manejo de cualquier dispositivo Android le resulte fácil y accesible.

Para conseguir esto mencionado, ha sido necesario modificar varias librerías que se encontraban obsoletas, tratando de utilizar otras más estándar y así obtener una interfaz más accesible. Un ejemplo de esto es el menú deslizante, que ha sido modificado por completo para obtener uno muy similar a los que nos podemos encontrar en aplicaciones como Google Maps o Gmail.

También se ha tratado de centrar el trabajo hacia la optimización en el uso de recursos y métodos dentro de la aplicación. Esto se verá más desarrollado en el anexo 4.

Para que el proyecto adquiera sentido al completo , ha sido necesario cooperar entre las dos partes durante todo el desarrollo del mismo. Se trata de una aplicación que envía peticiones y datos a un servidor, donde se procesa la información y se realizan los cálculos necesarios, por lo que la conexión entre ambos es importante.


\capitulo{2}{Objetivos del proyecto}


Objetivos de mejora:
\begin{itemize}
\item Actualización de la interfaz de proyecto desde un sistema Android obsoleto hacia la version mas actualizada posible del mismo, pudiendo implementar mejoras tanto visuales como funcionales
\item Actualizacion sobre la utilizacion de una galeria de imagenes relacionadas con los puntos de interes, ya que la ya implementada era utilizando el servicio de Panoramio, y habria que valorar si esos servicios siguen estando disponibles o si es necesario valerse de otros servicios, como Flickr o Instagram.
\item Diferenciación entre una ruta y otra si se ha realizado en la misma ciudad, para que la aplicacion calcule lugares diferentes basandose en los ya visitados.
\item Implementacion de aplicaciones externas como Twitter o TripAdvisor para poder descubrir las opiniones o valoraciones de otros usuarios.
\item Diferenciacion entre puntos de interes basandose en la actualizacion del algoritmo que pasa a utilizar horarios de apertura y cierre de los mismos, para que la aplicacion no visualice ningun lugar que no pueda ser visitado por estos factores.
\end{itemize}

Objetivos personales:
\begin{itemize}
\item Aprender a utilizar el entorno y lenguaje de programacion de Android, ya que es un sistema que puede resultar muy util para la vida laboral de un programador.
\item Mejorar en el proceso de la gestion de tareas, con aplicaciones como GitHub.
\item Aprender a utilizar el entorno de desarrollo de documentación LaTeX, ayundandome de la herramienta Texmaker.
\item Aprender a trabajar en modo SCRUM, utilizando Issues y Sprints para la realizacion del proyecto.
\item Aprender a realizar trabajo en equipo, manteniendo comunciacion con el compañero encargado de desarrollar el algoritmo del proyecto, para conseguir experiencia sobre un proyecto mas realista que si se realizara solo.
\end{itemize}


\capitulo{3}{Conceptos teóricos}

En este apartado se describirán ciertos conceptos teóricos que resultan necesarios para la correcta comprensión del proyecto.

\section{Modelo Cliente-Servidor}

Una parte importante de este proyecto es el concepto del modelo cliente-servidor, ya que cumple con la estructura que estamos utilizando. En este modelo y en el proyecto hay unas caracteristicas principales y unos actores concretos que son los siguientes:

\subsection{Cliente}
\begin{itemize}
\item Este se encarga de lanzar solicitudes o peticiones al servidor. Si lo comparamos con el proyecto de TourPlanner, podemos ver que la parte del cliente correspondería a la aplicación de Android.

\item El cliente es capaz de lanzar peticiones a diferentes servidores, como es nuestro caso en el que lanza peticiones tanto al servidor de OpenStreetMaps para poder visualizar el mapa online, a la base de datos que desarrolla mi compañero(lanza la petición a través de Glassfish) o por ejemplo a un servidor que provee las imágenes de los puntos de interés.

\item Obviamente este cliente contiene la interfaz que utilizará el usuario final, por lo que la interacción es directa.
\end{itemize}

\subsection{Servidor}
\begin{itemize}
\item El servidor es la otra parte indispensable de este modelo, ya que en el momento en que esta iniciado se encarga de escuchar a la espera de peticiones del cliente.

\item Al igual que un cliente puede realizar distintas peticiones a distintos servidores, el servidor puede mantener conexión con un gran numero de clientes y procesar las peticiones. 

\item Para este proyecto, quien cumple la funcion de servidor es GlassFish, que es a su vez un intermediario entre la aplicacion de android y la base de datos.
\end{itemize}

El proceso básico que también se verá en detalle mas adelante es el siguiente: el usuario interactua directamente con la aplicación (cliente), la cual realiza la petición a GlassFish (servidor), que a su vez realiza la petición a la base de datos y se la devuelve al cliente.

Al estar usando este modelo conseguimos una serie de ventajas:
\begin{itemize}
\item Centralización: Al contar con un servidor que recibe todas las peticiones que le lance el cliente, conseguimos un control de todos los accesos que se realizan y los recursos que se utilizan, pudiendo restringir o aumentar donde sea necesario.
\item Es un sistema muy fácilmente escalable, pudiendo aumentar las capacidades tanto de un lado como del otro independientemente, sin que al otro le afecte, o en todo caso teniendo que hacer modificaciones mínimas para que sea funcional. Esto mismo se traduce igualmente a la facilidad de mantenimiento, por las mismas razones.
\item Teniendo cierta relación con la primera ventaja, conseguimos seguridad, ya que actualmente contamos con tecnologías que permiten conexiones seguras entre cliente y servidor, y a su vez evitando posibles hackeos, ya que aunque varios clientes estén accediendo al mismo servidor, no pueden conocer las direcciones IP del resto. 
\end{itemize}

\section{Estándar}

Cuando se habla de estándar podemos considerar varias opciones, pero a la que nos referimos en este caso es el hecho de desarrollar software de un tipo especifico siguiendo ciertas reglas establecidas para así conseguir tener una cantidad indeterminada de programas y aplicaciones que contengan una estructura y una programación parecida, ya sea por el uso de los mismos métodos de una misma librería, o por seguir ciertos patrones con los que estructuras los programas de una manera concreta.

En el caso de Android, podemos considerar como estándar la librería en si que ofrece Android para programar, ya que todas las aplicaciones que están desarrolladas por Google para este tipo de dispositivos mantienen una estructura bastante parecida en el desarrollo de ciertas técnicas.

Una forma visual de demostrar esto se encuentra en el menú que utilizan la mayoría de aplicaciones, como el de Google Maps, Gmail o la aplicación para visualizar las fotos. Todos estos ejemplos utilizan una estructura prácticamente idéntica, y con esto se consigue lo que se llama estandarización, lo cual será explicado en el apartado de técnicas y herramientas.

\section{Code Smell}

Cuando se habla de Code Smell se esta hablando de problemas que puede contener el codigo de un programa, que no causa un error a la hora de ejecutar o compilar el programa, pero es un problema de caracter mas profundo, que acarrea otra serie de fallos, como puede ser la lentitud en un programa.

Estos problemas también dificultan el mantenimiento del código a futuro, ya que puede ser un error que en la actualidad no cause fallos, pero en el momento que se realicen modificaciones un programa que contenga estos Code Smells será mas propenso a fallos críticos, que, dependiendo del tipo de Code Smell que lo haya provocado, puede ser mucho mas difícil detectarlo en ese momento que cuando se cometió.

Gracias a la identificación de los Code Smells, se puede saber cuándo es necesario refactorizar código, que es una técnica que se explicará mas adelante. Con respecto a la identificación también es importante mencionar que tiene un carácter muy subjetivo, dependiendo del programador o incluso del lenguaje de programación, aunque existen herramientas que son capaces de detectar bastantes tipos de defecto, ya que algunos son mas subjetivos que otros (en el caso de código duplicado, aunque también pueden existir falsos negativos).

\section{API}

Una API, como su propio nombre indica, provee de una interfaz de programación de aplicaciones, en concreto, de una aplicación que será la que provee de esa API a otro software para que pueda hacer uso de rutinas, métodos y procedimientos con los cuales incluya la funcionalidad del software proveedor de la API.

Visto de otra forma, se basa en el uso de bibliotecas, y con esto se consigue que el programador no tenga que desarrollar todo el código que se encuentra en ellas, ya que con hacer uso de sus métodos y llamadas ya consigue esas funcionalidades.

Por otro lado, al hacer esto también se consigue estandarización, debido a que todos los programas que usen esas librerías para implementar sus funcionalidades estarán haciendo los mismos procesos de llamadas y usando los mismos métodos, así que se benefician de todas las ventajas con las que cuenta la estandarización, pero a su vez también corremos el riesgo de que esa estandarización se rompa, bien porque la API deje de dar soporte a la biblioteca que estamos utilizando, o porque en las nuevas versiones de una aplicación ya no sea compatible nuestra API implementada. 

Llevándolo a nuestro caso, tendríamos por un lado la API que nos permite visualizar imágenes de los puntos de interés, o por ejemplo la API con la que estamos viendo los mapas (OpenStreeMaps). Aquí se cumplen los dos casos referentes a la estandarización, con OpenStreetMaps conseguimos una API y un código válido y más sencillo, que se parecerá estructuralmente al de otros programadores y si necesitáramos soporte lo obtendríamos fácilmente, pero por el otro lado, tenemos la API de Panoramio que ya no es válida, dado que el servicio de esta página web dejó de estar disponible y por tanto nos veríamos obligados a buscar otra API que nos provea de las funcionalidades buscadas.

\section{Mantenimiento}

El concepto de mantenimiento en informática posee mucho valor, ya que el desarrollo de cualquier aplicación requiere mantenimiento posterior, y lo podemos dividir en dos aspectos:

\subsection{Soporte}

El soporte es una parte del mantenimiento, ya que después de cualquier desarrollo siempre habrá errores, bugs o fallos, y éstos pueden ser problemas que no afecten al funcionamiento normal de un aplicación o por el contrario que causen su parada total, quedando inutilizable. 

El soporte es un servicio a través del cual un especialista informático brinda asistencia, de forma física o remota a un individuo o una organización que hacen uso de la herramienta. En el caso de nuestra aplicación, el soporte del que hablamos sería remoto, a traves de una aplicación web o un foro, donde además de personas especializadas en esta aplicación, también hay otros usuarios que aportan su experiencia sobre ciertos problemas y en muchos casos estos problemas pueden ser resueltos sin necesidad de modificar el código, porque la falta de conocimiento del usuario es una de las mayores causas de incidencias que se suelen resolver en entornos de soporte.

\subsection{Actualización}

La otra parte importante del mantenimiento de una aplicación, sobre la que haremos mención en muchos aspectos de este trabajo es la actualización. 

Se basa en el cambio de código, APIs y librerías para obtener métodos nuevos, que están más optimizados para realizar la misma funcionalidad o también para solucionar errores que existieran en versiones anteriores. Por último, también se aplica a la implantación de nuevas funcionalidades de una misma aplicación, para que cada vez cumpa mejor la labor para la que haya sido desarrollada. En conjunto, supone una parte fundamental y necesaria para la aplicación. Cuando la actualización no se realiza con cierta periodicidad la aplicación acaba sufriendo pequeños errores o errores críticos. Cuanto más largo sea el periodo entre una actualización y la siguiente más complicado será conseguir que todo lo que se haya ido cambiando o mejorando sea aplicado de manera fácil.



\section{Referencias}

Las referencias se incluyen en el texto usando cite \cite{wiki:latex}. Para citar webs, artículos o libros \cite{koza92}.


\section{Imágenes}

Se pueden incluir imágenes con los comandos standard de \LaTeX, pero esta plantilla dispone de comandos propios como por ejemplo el siguiente:

\imagen{escudoInfor}{Autómata para una expresión vacía}



\section{Listas de items}

Existen tres posibilidades:

\begin{itemize}
	\item primer item.
	\item segundo item.
\end{itemize}

\begin{enumerate}
	\item primer item.
	\item segundo item.
\end{enumerate}

\begin{description}
	\item[Primer item] más información sobre el primer item.
	\item[Segundo item] más información sobre el segundo item.
\end{description}
	
\begin{itemize}
\item 
\end{itemize}

\section{Tablas}

Igualmente se pueden usar los comandos específicos de \LaTeX o bien usar alguno de los comandos de la plantilla.

\tablaSmall{Herramientas y tecnologías utilizadas en cada parte del proyecto}{l c c c c}{herramientasportipodeuso}
{ \multicolumn{1}{l}{Herramientas} & App AngularJS & API REST & BD & Memoria \\}{ 
HTML5 & X & & &\\
CSS3 & X & & &\\
BOOTSTRAP & X & & &\\
JavaScript & X & & &\\
AngularJS & X & & &\\
Bower & X & & &\\
PHP & & X & &\\
Karma + Jasmine & X & & &\\
Slim framework & & X & &\\
Idiorm & & X & &\\
Composer & & X & &\\
JSON & X & X & &\\
PhpStorm & X & X & &\\
MySQL & & & X &\\
PhpMyAdmin & & & X &\\
Git + BitBucket & X & X & X & X\\
Mik\TeX{} & & & & X\\
\TeX{}Maker & & & & X\\
Astah & & & & X\\
Balsamiq Mockups & X & & &\\
VersionOne & X & X & X & X\\
} 

\capitulo{4}{Técnicas y herramientas}

Esta parte de la memoria tiene como objetivo presentar las técnicas metodológicas y las herramientas de desarrollo que se han utilizado para llevar a cabo el proyecto. Si se han estudiado diferentes alternativas de metodologías, herramientas, bibliotecas se puede hacer un resumen de los aspectos más destacados de cada alternativa, incluyendo comparativas entre las distintas opciones y una justificación de las elecciones realizadas. 
No se pretende que este apartado se convierta en un capítulo de un libro dedicado a cada una de las alternativas, sino comentar los aspectos más destacados de cada opción, con un repaso somero a los fundamentos esenciales y referencias bibliográficas para que el lector pueda ampliar su conocimiento sobre el tema.


En este apartado se destacarán las principales técnicas y herramientas utilizadas para el proyecto. Cabe destacar que habrá herramientas también utilizadas que no serán explicadas, bien por pertenecer a la parte del servidor o por ser herramientas que fueron utilizadas en la anterior versión de este trabajo, pero que para esta versión sus funcionalidades no se han modificado y por tanto no se ha aportado nada adicional a lo que ya se implementó.

\subsection{Técnicas}

Las técnicas utilizadas durante el desarrollo de un proyecto pueden ayudar a comprender mejor la metodología que se ha llevado a cabo, así como la robustez que tiene el código. También nos permiten observar ligeramente el rumbo que ha tomado el proyecto, ya que puede ser un proyecto mas enfocado a la implementación de métodos desde cero, o para complementar un proyecto con un ciclo de vida mas largo. Por otro lado, también se puede estar buscando el mantenimiento, optimización y actualización del código que ya esta desarrollado en un proyecto, como es nuestro caso.

\subsubsection{Refactorización}

La refactorización es una técnica utilizada en programación para realizar modificaciones sobre el código de un programa en cuanto a la estructura interna, pero sin que el comportamiento externo se vea afectado. Son cambios que se pueden apreciar estructuralmente y que a simple vista podrían no tener importancia, pero a medida que un programa avanza, se actualiza o se dejan de usar ciertas técnicas para pasar a usar otras, haber hecho esta refactorización cobra mas importancia. 

Si nos paramos a pensar, el termino de refactorización tiene el origen en las matemáticas, ya que la "factorización" en las matemáticas lo que puede hacer es convertir un polinomio compuesto por una variable y un numero en 2 polinomio mas simples y con una estructura mas clara de ver. El funcionamiento externo es en sí el mismo, pero se consigue una limpieza o una robustez estructural. Esto es lo que se busca al hacerlo en la programación.

La refactorización abarca desde cambiar el nombre de una variable para que vaya mas acorde con la función que llevará a cabo, hasta mover un método de una clase a otra para así conseguir un acceso mas optimizado y una estructura mas lógica.

Hay una serie de refactorizaciones sencillas pero que son de las que mas se realizan, debido a que solucionan los code smells que más se cometen:
\begin{itemize}
\item EXTRACT METHOD: Es el proceso por el cual, tras detectar que ciertas porciones de código se repiten en distintos métodos o incluso dentro de un mismo método en partes distintas del mismo, pero que tienen la misma estructura y en definitiva cumplen la mismo función, por lo que esas porciones de código repetido pueden ser extraídas para crear un método nuevo, de tal manera que cuando se necesite hacer el proceso que exija usar ese mismo proceso, bastará con llamar al método. 

Así se consigue una estructura más robusta en el código, así como mejor mantenibilidad del mismo, ya que si en un momento dado es necesario modificar esas lineas, en vez de modificarlas en cada una de las localizaciones donde estuvieran, solo se hará en el método nuevo.

\item MOVE METHOD: En este caso, durante la detección de los posibles problemas que pueda tener el código, se observa que un método que se encuentra en una clase, no esta correctamente implementado en la misma, ya que a lo mejor esta haciendo mas uso de recursos de otra clase distinta, y así se realiza, se decide mover ese método a la clase en la que debería estar implementado, ya que esto hace que la estructura sea más lógica, y por consiguiente mas óptima.

\item RENAME: Este tipo de refactorización es el más sencillo y a la vez es el que más necesario resulta en muchas ocasiones, ya que se basa en renombrar una variable, método o clase para que tenga más sentido con respecto a la labor que desempeñe. El hecho de que existan nombres excesivamente genéricos o que directamente no cumplan con lo que es realmente es un error muy frecuente, por tanto, esta refactorización es muy frecuente.
\end{itemize}

\subsubsection{Migración}

El concepto de migración puede ser muy extenso, pero en este caso trataremos acerca de la migración en el software.

Este tipo de migración se realiza cuando se quiere cambiar por completo un software para pasar a utilizar otro que este mas actualizado, o cumpla mejor con las necesidades que exija el proyecto en cuestión. 

Para relacionar la técnica con este proyecto, hay que mencionar Android, ya que esta técnica se ha aplicado para realizar la migración desde Android hacia AndroidX, ya que éste último es la version mejorada del anterior. Pero no es una simple actualización, ya que se encuentran en librerías diferentes, pero esto se explicará mas adelante en el apartado de herramientas.

\subsubsection{Actualización}

Con respecto a este concepto entendemos actualización como la tecnica por la cual un software obtiene caracteristicas nuevas o arregla algunas que ya existieran a traves de las distintas versiones. Normalmente las actualizaciones resultan ser importantes ya que suelen corregir errores que pudiera tener el codigo anterior o mejoran alguas funcionalidades ya existentes.

Hay ocasiones en que las actualizaciones, aunque necesarias, no son obligatorias, y por tanto se puede seguir funcionando sin aplicarlas. Otras por el contrario, obligan a aplicarse, ya sea por ser razones de seguridad o porque en el momento que cambia un sistema tiene que estar con las ultimas versiones, porque si no suele causar incompatibilidades.

\subsection{Herramientas}

Como ya se ha comentado antes sobre las tecnicas y herramientas utilizadas, solo se destacara las que han tenido una verdadera relevancia para este proyecto, ya que hay algunas que ya se mencionan en la version anterior de TourPlanner.

\subsubsection{Android}

En esta seccion se hablará de Android como herramienta principal, pero tambien de dos herramientas directamente relacionadas, que son Android Studio y AndroidX.

Android Studio es un 





\capitulo{5}{Aspectos relevantes del desarrollo del proyecto}

En este apartado se van a destacar los aspectos que más relevancia han adquirido a lo largo del desarrollo del proyecto, bien por haber realizado cambios importantes o por haber supuesto un reto más en el trabajo.

Se verán ciertos aspectos que estarán más desarrollados en los anexos, sobre todo en el anexo 4. También se tratará de justificar las modificaciones que se hayan realizado sobre el proyecto anterior \cite{tfm1}, para que quede claramente diferenciada la aportación que ha supuesto este proyecto.

\section{Trabajo sobre un proyecto ya desarrollado}

En esta sección se quiere mostrar los aspectos más importantes a la hora de trabajar sobre otro proyecto que ya ha sido desarrollado, ya que pueden suponer ciertas dificultades si no se tienen en cuenta.

Para este proyecto, se ha comenzado con dos preámbulos ya mencionados, el primer desarrollo \cite{tfg1} donde se comenzaba el desarrollo de la aplicación y el segundo \cite{tfm1}, donde se mejoraba este desarrollo anterior. A la hora de plantear esta tercera versión del proyecto, ha sido necesario realizar un trabajo inicial de documentación para poder entrar en contexto sobre lo que se había desarrollado y lo que quedaba por desarrollar. Para ello, los apartados sobre objetivos y lineas de trabajo futuras de ambos proyectos han sido bastante útiles, ya que a través de estos brindaban la visión que tenían sobre sus proyectos y las mejoras que se podían desarrollar.

También ha sido necesario investigar la antigüedad de estos proyectos, concretamente de la segunda versión \cite{tfm1}, ya que era la última, aunque más adelante se vió que ésta había sido desarrollada sobre la misma versión tanto para el cliente como para el servidor. Por tanto, una parte crucial es la de actualizar, sobre todo en lo referente al cliente, ya que la versión de la que se habla es Android 4.4 (API 14), lo cual se sitúa 7 años atrás. Además durante estos últimos años, los dispositivos móviles han sufrido muchos avances tecnológicos y esto se ha podido ver incluso durante el desarrollo del proyecto, ya que cuando se comenzó la última version de Android era la de Android 9.0 y en la actualidad ya se ha lanzado Android 10. En este artículo \cite{and1} se puede ver de una manera sencilla la evolución que ha sufrido Android a lo largo de los años.

Uno de los cambios que ha supuesto toda esta evolución de Android surgió un año después del lanzamiento de Android 4.4, ya que en Android 5.0 se incluye el concepto de \textit{Material Design} que sirve desde entonces como base de diseño de aplicaciones de Android para unificar todos los dispositivos que se encontraran por encima de esta versión. Es un paso más hacia la estandarización que en este proyecto también se ha incluido.

Otro aspecto importante a la hora de trabajar sobre un proyecto anterior, es la necesidad de comprender cómo se ha desarrollado el mismo y para ello toman mucha relevancia los comentarios que se realizaron sobre el código, para evitar confusiones o simplemente para conseguir que la experiencia de mejorar esta aplicación sea más llevadera.

\imagen{comment1}{Ejemplo de comentarios sobre el código que ayudan a la comprensión del mismo}


\section{Actualización y cambio de librerías}

Uno de los aspectos más importantes para el desarrollo de este proyecto ha sido la actualización de varias librerías, ya que se encontraban obsoletas por haber dejado de recibir soporte o por estar desactualizadas. Tras haber visto el apartado sobre las librerías en los conceptos teóricos, se puede ver la importancia que poseen en un proyecto, por lo que sin éstas funcionando, no se puede desarrollar el resto de la aplicación.

Ha sido necesario realizar un trabajo de investigación sobre cada una de estas librerías obsoletas con el fin de encontrar una alternativa más estándar y actual, tratando de conseguir más accesibilidad y soporte en un futuro.

Llevándolo a este caso, se tendría por un lado una API con la que se obtienen los mapas (OpenStreeMaps). Aquí se cumplen los dos casos referentes a la estandarización, con OpenStreetMaps se consigue una API y un código válido y más sencillo, que se parecerá estructuralmente al de otros programadores y si se necesitara soporte lo obtendríamos fácilmente, pero por el otro lado, la API de Panoramio que ya no es válida, dado que el servicio de esta página web dejó de estar disponible y por tanto sería necesario buscar otra API que ofrezca las funcionalidades buscadas.

Las librerías más importantes que han tenido que ser sustituidas son:

\begin{itemize}
\item \textbf{SlidingMenu:} Librería que permitía implementar un menú deslizante para poder navegar entre ventanas.
\item \textbf{SherlockActivity:} Librería utilizada para la implementación de cada una de las ventanas (Activities) de la aplicación.
\item \textbf{Org.Apache:} Librería que se utilizaba para realizar todas las peticiones HTTP y HTTPS hacia el servidor.
\end{itemize}

Toda la explicación sobre el cambio de estas librerías se encuentra en el anexo 4: Manual del programador.

La sustitución de librerías al completo dentro de una aplicación supone un cambio en todos los métodos que se aprovechen de ellas, ya sea a través de llamadas o variables de las mismas, para pasar a utilizar las nuevas. El problema principal es que en muchas ocasiones, aunque se consiga una nueva librería, puede que hayan cambiado la forma de hacer las llamadas a métodos o los parámetros dentro de esas llamadas y es necesario tenerlo en cuenta para poder hacer un cambio que realmente tenga un efecto positivo.

\section{Cambio de la herramienta de desarrollo}

Para este proyecto, también se ha modificado la herramienta de desarrollo que se estaba utilizando, ya que anteriormente se utilizó Eclipse, por ser el entorno oficial de desarrollo de aplicaciones Android, pero en 2015 se lanzó la primera versión de Android Studio, que lo sustituyó como herramienta oficial.

Al iniciar este proyecto se pudo ver que se había utilizado Eclipse al ser 2014 y no existir Android Studio, pero ahora ya es necesario utulizar esta nueva herramienta, que también ofrece nuevas funcionalidades, y un entorno de desarrollo totalmente especializado para las aplicaciones de Android. 

Uno de los grandes cambios que supone esta herramienta, es \textit{Gradle} \cite{gradle2} que es un sistema de automatización de construcción de código abierto, cuyos plugins iniciales están principalmente centrado en desarrollo y despliegue en Java, Groovy y Scala \cite{gradle1}. En concreto, Android Studio es una de las aplicaciones que lo integran para el desarrollo de aplicaciones, ya que permite automatizar la gestión de dependencias sobre librerías, APIs y repositorios.

Se puede ver con más detalle el cambio que ha supuesto la utilización de Android Studio en el anexo 4.

\section{Programación en Android}

Para poder desarrollar este proyecto, ha sido necesario partir desde cero en cuanto a la programación con Android. Hay una serie de conceptos importantes sobre los que ha habido que aprender para poder realizar esta programación:

\begin{itemize}
\item \textbf{Lenguajes de programación:} Para desarrollar este proyecto ha sido necesario tener conocimientos sobre Java para implementar todo el código referente al funcionamiento interno, así como CSS para la parte visual.
\item \textbf{Actividades:} El concepto de actividades para Android, como ya se ha visto en los conceptos teóricos, es algo crucial al desarrollar una aplicación, porque representan cada una de las ventanas que aparecen al realizar interacción con la aplicación.
\item \textbf{Notificaciones:} El sistema que utiliza Android también es algo bastante importante, ya que hay que conocer la forma en la que las aplicaciones activan estas notificaciones y en base a que.
\item \textbf{Permisos:} Al igual que las notificaciones, pero en este caso es algo necesario, ya que en función a lo que se quiera mostrar u obtener de la aplicación, puede que necesitemos permisos para acceder a los distintos servicios del dispositivo. Estos servicios podrían ser desde la cámara hasta la geolocalización. Si no se establece una manera correcta para gestionar estos permisos, la aplicación que se desarrolle no funcionará como debería.
\end{itemize}

\section{Cambio de interfaz}

Ya que para esta aplicación ha sido necesario actualizar gran parte de las librerías y repositorios que se utilizaban, también se ha decidido realizar un cambio sobre la interfaz de la misma, para que concuerde mejor con una interfaz estándar de una aplicación Android, tanto en los colores como en el menú. Cada una de las ventanas mantiene un aspecto similar en cuanto a la estructura de lo que contenían, porque también se consideró que no era necesario que pareciera una aplicación totalmente diferente a la versión anterior. 

Como ya se ha mencionado previamente, a partir de la version de Android 5.0, surgió \textit{Material Design}, que trataba de aportar una manera de desarrollar la interfaz de las aplicaciones más estándar. Esto se ha querido aplicar para este proyecto, ya que se está basando sobre una versión mínima de Android 6.0, pudiendo funcionar hasta en la última versión disponible.

Este cambio sobre la interfaz se puede ver en el anexo 4 y también en el anexo 5.

\section{Trabajo colaborativo}
 
Otro aspecto relevante de este proyecto ha sido el trabajo colaborativo, que se ha realizado para desarrollar el proyecto complementario sobre el BackEnd \cite{tfg2} y este proyecto enfocado en el FrontEnd.

Dado que el cliente y el servidor tienen que tener una conexión, también es necesario que el desarrollo sobre los mismos la tenga, por lo que ha sido necesario coincidir en varios puntos, incluso mantener una comunicación continua al término final del proyecto, para contrastar ciertas partes de la memoria y para brindarse ayuda en el desarrollo de partes que influían en la parte complementaria.

Para que se comprenda mejor el sentido de esta colaboración, el proceso básico que también se verá en detalle mas adelante es el siguiente: el usuario interactúa directamente con la aplicación (cliente), la cual realiza la petición a GlassFish (servidor), que a su vez realiza la petición a la base de datos y se la devuelve al cliente. 

Una de las partes que más colaboración ha tenido para este proyecto, ha sido la modificación de la interfaz para que ésta ofreciera nuevos métodos de cálculo sobre las rutas turísticas, pudiendo utilizar los nuevos algoritmos desarrollados en el proyecto de BackEnd \cite{tfg2}.

A la hora de realizar comprobaciones en la conexión de ambas partes, ha sido necesario también mantener una comunicación, ya que a la hora de lanzar peticiones, el cliente simplemente queda a la espera de que el servidor pueda devolvérselos, pero en el momento de obtener fallos de conexión, ya es necesario observar lo que ha podido ocurrir tanto en un extremo como en el otro. Esto se consigue a través de la colaboración.

Cabe destacar que en el desarrollo tanto de la memoria como de los anexos se hacen referencias continuas hacia el otro proyecto, incluso algunos apartados contienen la misma información al tratarse de un desarrollo por separado pero sobre un mismo objetivo, que es el de que funcione la aplicación que se desarrolla con todas sus funcionalidades.
\capitulo{6}{Trabajos relacionados}

Este apartado sería parecido a un estado del arte de una tesis o tesina. En un trabajo final grado no parece obligada su presencia, aunque se puede dejar a juicio del tutor el incluir un pequeño resumen comentado de los trabajos y proyectos ya realizados en el campo del proyecto en curso. 

\capitulo{7}{Conclusiones y Líneas de trabajo futuras}

Todo proyecto debe incluir las conclusiones que se derivan de su desarrollo. Éstas pueden ser de diferente índole, dependiendo de la tipología del proyecto, pero normalmente van a estar presentes un conjunto de conclusiones relacionadas con los resultados del proyecto y un conjunto de conclusiones técnicas. 
Además, resulta muy útil realizar un informe crítico indicando cómo se puede mejorar el proyecto, o cómo se puede continuar trabajando en la línea del proyecto realizado. 



\bibliographystyle{plain}
\bibliography{bibliografia}

\end{document}
